%     Tokyo Debian Meeting resources
%     Copyright (C) 2012 Junichi Uekawa
%     Copyright (C) 2011, 2015, 2020 Nobuhiro Iwamatsu

%     This program is free software; you can redistribute it and/or modify
%     it under the terms of the GNU General Public License as published by
%     the Free Software Foundation; either version 2 of the License, or
%     (at your option) any later version.

%     This program is distributed in the hope that it will be useful,
%     but WITHOUT ANY WARRANTY; without even the implied warranty of
%     MERCHANTABILITY or FITNESS FOR A PARTICULAR PURPOSE.  See the
%     GNU General Public License for more details.

%     You should have received a copy of the GNU General Public License
%     along with this program; if not, write to the Free Software
%     Foundation, Inc., 51 Franklin St, Fifth Floor, Boston, MA  02110-1301 USA

%  preview (async-shell-command (concat "evince " (replace-regexp-in-string "tex$" "pdf"(buffer-file-name))) "*tex-preview*")

%%ここからヘッダ開始。

\documentclass[mingoth,a4paper]{jsarticle}
\usepackage{monthlyreport}
% 日付を定義する、毎月変わります。
\newcommand{\debmtgyear}{2021}
\newcommand{\debmtgmonth}{5}
\newcommand{\debmtgdate}{15}
% started from zero:
% (let ((year 2021) (month 5)) (+ (* (- year 2005) 12) month -1))
\newcommand{\debmtgnumber}{196}

% Needed to import pandoc-generated LaTeX documents.
% See https://stackoverflow.com/questions/40438037/tightlist-error-using-pandoc-with-markdown
\providecommand{\tightlist}{%
  \setlength{\itemsep}{0pt}\setlength{\parskip}{0pt}}

\begin{document}

\begin{titlepage}
\thispagestyle{empty}
% タイトルページ:編集必要な部分は最初のマクロに飛ばすこと

\vspace*{-2cm}
第\debmtgnumber{}回 東京エリア Debian 勉強会資料\\
\hspace*{-2cm}
\includegraphics{image-assets/dotdeb.pdf}\\
\hfill{}\debmtgyear{}年\debmtgmonth{}月\debmtgdate{}日

% ここはアップデートすること
% 全角文字にしないとフォントのサイズが合わないので注意
\rotatebox{10}{\fontsize{30}{30} {\gt goを使う}}\\

\vspace*{-2cm}
\hfill{}\includegraphics[height=6cm]{image-assets/openlogo-nd.eps}
\end{titlepage}

\newpage

\begin{minipage}[b]{0.2\hsize}
 \definecolor{titleback}{gray}{0.9}
 \colorbox{titleback}{\rotatebox{90}{\fontsize{80}{80} {\gt デビアン勉強会} }}
\end{minipage}
\begin{minipage}[b]{0.8\hsize}
\hrule
\vspace{2mm}
\hrule
\begin{multicols}{2}
\tableofcontents
\end{multicols}
\vspace{2mm}
\hrule
\end{minipage}

\dancersection{最近のDebian関連のミーティング報告}{杉本　典充}

\dancersection{事前課題}{杉本　典充}

今回の事前課題は以下です。

\begin{enumerate}
 \item go 書いたことありますか、何を書きましたか
 \item go で書かれたプログラムでDebianにあったらいいなというのはなにかありますか？
\end{enumerate}

%この課題に対して提出いただいた内容は以下です。

\begin{multicols}{2}
{\small
% \begin{prework}{ ����(yy\_y\_ja\_jp) }
���ޤ���������ǤϤʤ��Ǥ���... �긵�Υ�åץȥåפǤ� /boot �ˤ� ext2
 ����¾�ˤ� ext3���Ƕ�ȤäƤ���ǥ����ȥåפǤ� /boot �ˤ� ext3 ����¾
 �ˤ� LVM ��� ext4 ��ȤäƤ��ޤ���
\end{prework}

\begin{prework}{ �����ϥ� }
�ǥե���Ȥ�ext3����
�����ǥե���ȤΤޤޡ�
(NTFS���VM���᡼�����ext3�⤢�뤱��)
\end{prework}

\begin{prework}{ yos.takahashi }
ext3/4���˻ȤäƤޤ���ext3�Υǡ�������١����ˤĤ�������Linux2011ǯ1���˼�ɮ���ޤ�����
\end{prework}

\begin{prework}{ MATOHARA }
����inode �ϳ���������Ƥ���inode ��ưŪ�˳�����Ƥ���XFS �����򤹤뤳
 �Ȥ�¿���Ǥ���NILFS �Ͼ�����Ƥߤ��ΤǤ�����mount ���˰ʲ��Τ褦�ʥ��
 ���������ФƤޤ��ݤ��ʤȻפ��ޤ�����
\begin{commandline}
$ sudo mount /dev/sdb1 /mnt
mount.nilfs2: WARNING! - The NILFS on-disk format may change at any time.
mount.nilfs2: WARNING! - Do not place critical data on a NILFS filesystem. 
\end{commandline}
����¾NotePC �Ǥ�dm-crypt �ξ�˥ե����륷���ƥ���֤��ưŹ沽�����ꡢ
 eCryptfs �ǰŹ沽�����ꤷ�Ƥ��ޤ����񤭹��߻���CPU �򤫤ʤ���񤷤ޤ��ġ�
\end{prework}

\begin{prework}{ ��ޤ� }
��ext2-$>$reiserfs-$>$jfs-$>$xfs-$>$reiserfs-$>$ext3�ȻȤäƤ��ޤ�����

����:
\begin{itemize}
 \item reiserfs: ���ե������¿���ե�����Υ������������ӥ��Ӥ��Ƥ��ɤ���
       �����������դȤ����ⵤ�������θ�μ�žȬ�ݥ����ɤˡ�����
 \item jfs: ����v1.0��̾��ä�IBM���ꥨ�ʥ��������ˤ�xfs��ƨ��
 \item xfs: fsck==true�˴�ư����������ǯ�Ȥä���ΤΥޥ�����Ĵ����0byte
       �ե��������������Ѥ���줺ƨ����
\end{itemize}
������reiser4��Ķ���Ԥ��뤦���ˤ��줬�����ʤäơ����ext3�˸��경��htree�����ä������⤦Ŵ�Ĥʤ鲿�Ǥ⤤���Ǥ��������Ȥ����Ĥ�nilfs�ʤɤ˼��Ф��Ƥ��ޤ���ext3�������noatime���٤Ǥ����������aufs���碌�Ƽ�ʬ���Ȥ��Ȥ� *strap �Ķ��򥯥����˥󥰤��Ƥ��ޤ����¸����ƥ��Ȥ������Ǥ�����Ǥ���USB�����ư�Ǥ�ͭ�ѡ�

���LVM�ǤϤʤ�MD��Ȥäƾ�Ĺ���ܥХå����åפ򤷤Ƥ��ޤ���
 MD(sda,sdb,sdc)�ǹ��ۤ����̾��MD(����)�Dz�ư���Хå����åפλ���
 attach/detach�򤹤롣�֥��å���٥�ʤΤǥꥫ�Х��FSǤ���Ǥ��������̤�
 �������¾����ˡ���ʤ�������
\end{prework}

\begin{prework}{ henrich }
�ȤäƤ���֤�NTFS��Ĺ���󤸤�ʤ��Ǥ����͡����졣
�����Ρ����Ѥο������ǥ�������ext4�ǥե����ޥåȤ��ޤ����������㤤��Ƚ��ʤ��Ȥ����������Ƥ��ޤ���
\end{prework}

\begin{prework}{ emasaka }
�Ĥ뤷��FS��ȤäƤޤ�
\end{prework}

\begin{prework}{ �ܾ� }
ext3����Ѥ��Ƥ��ޤ����ä��Ѥ�ä����ȤϤ��Ƥ��ޤ���
FS����ʤ��Ǥ������Ƕ�Lenny��2TB��HDD��Ȥä���parted�äƤλȤ��ƶä��ޤ�����
\end{prework}

\begin{prework}{ ����@������ }
���������Ū�˳��Ѥ��Ƥ���ե����륷���ƥ��ReiserFS�Ǥ���
�Ż��ǻȤäƤ�Ķ���ext3�Ǥ�����ext3���Ÿ������ǥ��㡼�ʥ뤬
����Ʋ��Ǥ������ηи��ʸŤ������ͥ�Ǥ���...�ˤ����ꡢ
���ޤ꿮�Ѥ��Ƥ��ޤ���
�����ReiserFS�Ķ��ǤϤ���ޤǤνꤤ���ʤ��Ÿ����ڤä��ꡢ
�Ƥ�HDD�����줫�����ꤷ�Ƥ��ﳲ����ä��и���̵���Τǡ�
��³Ū�˻ȤäƤ��ޤ���
��ǯ����ReiserFS�Υᥤ��ȯ��(Hans Reiser)�����ᤵ��Ƥ��ޤ������ƥʥ󥹤��ۤ��Ƥ��ޤ�����
�����������θ��ReiserFS��¾�γ�ȯ�Ԥˤ���³�����ݼ餵��Ƥ���Τǡ��¿����ޤ�����
\end{prework}

\begin{prework}{ nozzy123nozzy }
\begin{enumerate}
 \item LVM�ˤĤ��Ƥϡ�CentOS5.5��Ƴ�������Τ��Τޤޤ����Ѥ��Ƥޤ�����
       ���������ƥब��äƤ���Volume̾�ϥǥե���Ȥ�����ȡʼºݤˤ�
       kickstart�ˤơˤ��ѹ����ƻȤäƤޤ����ʾ㳲���Υ���١����˺��뤿
       ���
\item ext3�ˤĤ��Ƥϡ�debian-sid�����Τޤ޻��ꤷ�Ƥ����Τ򤽤Τޤ޻Ȥ�
      �Ƥ����ꤷ�ޤ���relatime, noatime ���餤�Ͼ����ɲä��Ƥߤ����ʡ���
      �ϻפäƤޤ���
\end{enumerate}
\end{prework}

\begin{prework}{ �ޤ��������ؤ� }
\begin{itemize}
 \item Debian�Ǥ��ä˶Ťä����Ȥ�����ext3��ȤäƤޤ������ۥޥ����qcow2
       ���᡼���ǥ��������ѻ��ʳ��ϡ�LVM�ϻȤäƤޤ���
 \item �����ǰ����ü�ʤΤϡ��������DHCP�������Ѥ�Armadillo-J�ǻȤäƤ�
       ��JFFS�Ǥ����ǥե���ȤΥե����०�����Ǥϥ�֡��Ȥ�����������
       �ƽ��������Ƥ��ޤ��Τǡ�RAM�ΰ�˽񤭤��ߡ��Ÿ��ڤäƤ�ä��ʤ�
       ���������Ǥ���Debian��udhcp�Υ������ѥå���������ӥ�ɤ��ƻȤäƤޤ���
       \footnote{\url{http://d.hatena.ne.jp/mkouhei/20080601/1212330630}}
 \item ��Debian���ߤǡ���ʬ����ǰ��֥ۥåȤʤΤ�palm webOS�Ǥ��������Ubuntu��
       �������ޥ���������Τ餷���ΤǤ�����/etc/mtab�򸫤��35�Ԥ⤢�ꡢ
       ���ʤ����֤ʹ����Ǥ��͡�
\end{itemize}
\end{prework}
}
\end{multicols}

\dancersection{Debian Trivia Quiz}{username}

Debianの昨今の話題についてのQuizです。

今回の出題範囲は\url{debian-devel-announce@lists.debian.org} や \url{debian-news@lists.debian.org}などに投稿された内容からです。

\begin{multicols}{2}
 %; whizzy-master ../debianmeetingresume201211.tex
% $B0J>e$N@_Dj$r$7$F$$$k$?$a!"$3$N%U%!%$%k$G(B M-x whizzytex $B$9$k$H!"(Bwhizzytex$B$,MxMQ$G$-$^$9!#(B
%

\santaku
{DebConf13 $B$N3+:ECO$H3+:EF|$O!)(B}
{$BF|K\!"El5~ET(B 6$B7n(B20$BF|(B}
{$B%K%+%i%0%"(B $B%^%J%0%"(B 7$B7n(B8-14$BF|(B}
{$B%9%$%9!"%t%)!<%^%k%-%e(B 8$B7n(B11-18$BF|(B}
{3}
{$B%K%+%i%0%"$O(BDebConf12$B$N3+:ECO$G$9!#(B
DebConf13$B$O%9%$%9$N%-%c%s%WCO$G3+:E$G$9!#(B
6/20$B$O3'$5$sM=Dj$r6u$1$F$*$-$^$7$g$&!#(B}

\santaku
{$B@$3&$N(BWeb$B%5!<%P$G:G$b?M5$$N$"$k(BLinux $B%G%#%9%H%j%S%e!<%7%g%s(B(W3Techs$BD4$Y(B)$B$O!)(B}
{CentOS}
{Debian}
{Ubuntu}
{B}
{\url{http://w3techs.com/technologies/history_details/os-linux}$B$K7k2L$N%0%i%U$,$"$j$^$9!#(B
$B8=:_(B Linux $B$r;HMQ$7$F$$$k(B web $B%5!<%P$N(B 32.9\% $B$,(B Debian $B$rMxMQ$7$F$*$j!"$=$N3d9g$O8=:_$bA}2C$rB3$1$F$$$k$=$&$G$9!#(B}

\santaku
{Debian $B%+!<%M%k%A!<%`$N%a%s%P!<$G$"$j!"(Bkernel.org $B$N(B 3.2.y $B0BDjHG7ONs$N%a%s%F%J$G$b$"$k(B Ben Hutchings $B$5$s$,<!4|(B Debian $B0BDjHG$H0l=o$K=P2Y$5$l$k(B Linux $B%+!<%M%k$K(B (3.2 $B7ONs$N(B mainline $B$K$OL5$$(B) $BDI2C5!G=$,Ek:\$5$l$kM=Dj$G$"$k$H=R$Y$F$$$^$9!#(B
$BB?$/$NDI2CE@$NCf$K4^$^$l$J$$$b$N$O2?!)(B}
{PREEMPT\_RT}
{Hyper-V guest drivers$B$N6/2=(B}
{ARM64/AArch64$B%"!<%-%F%/%A%c%5%]!<%H(B}
{C}
{Hyper-V guest drivers$B$O(Bmainline kernel$B$G(B3.2$B$K$b4^$^$l$F$$$^$9$,!"$h$j2~A1$5$l$?(B3.4$B$+$i$N=$@5$,F3F~$5$l$^$9!#(B
PREEMPT\_RT$B$O%O!<%I%j%"%k%?%$%`$r<B8=$9$k$?$a$N(BPatch$B!"(B
linux-image-rt-amd64 , linux-image-rt-686-pae $B$N(Bmetapackage$B$G;HMQ$G$-$^$9!#(B
$B?7$7$$(BARM 64$B%S%C%H%"!<%-%F%/%A%c%5%]!<%H$O(Bmainline kernel 3.7$B$+$i(B}

\santaku
{Wookey$B$5$s$,%"%J%&%s%9$7$?(Balpha$BHG$N(BDebian port arm64 image$B$O!)(B}
{Debian/Ubuntu port image}
{Debian/KFreeBSD port image}
{Debian/GnuHurd port image}
{A}
{self-bootstrapp(non x86)$BBP1~$H$N$3$H$G$9!#(B\url{http://wiki.debian.org/Arm64Port}$B$G%9%F!<%?%9$,3NG'$G$-$^$9!#(B}

\santaku
{700,000$BHVL\$N%P%0$,Js9p$5$l$?F|$rEv$F$k(B700000thBugContest$B$N7k2L$,=P$^$7$?!#$=$NM=A[F|$HJs9pF|$O!)(B}
{2012/12/12$B$rM=A[$7$?(BDavidPrevot}
{$BM=A[F|(B:2013/02/04$B!"Js9pF|(B:2013/02/14}
{$BM=A[F|(B:2013/02/07$B!"Js9pF|(B:2013/02/14}
{$BM=A[F|(B:2013/02/14$B!"Js9pF|(B:2013/02/07}
{C}
{$B:G$b6a$$(B2013/02/14$B$rM=A[$7$?(BChristian Perrier$B$5$s$,Ev$F$^$7$?!#7k2L$O(B\url{http://wiki.debian.org/700000thBugContest}$B$G8x3+$5$l$F$$$^$9!#(B
$B$^$?!"(B800,000/1,000,000$BHVL\$N%P%0$,Js9p$5$l$kF|$rEv$F$k%3%s%F%9%H(B\url{http://wiki.debian.org/800000thBugContest}$B$b3+:E$5$l$F$$$^$9!#(B}

\santaku
{master.debian.org$B$,?7$7$$5!3#$K0\9T$5$l$^$7$?!#$3$l$O2?$N%5!<%P$G$7$g$&$+(B $B!)(B}
{@debian.org$B$N%a!<%k%5!<%P(B}
{$B%Q%C%1!<%8$N%^%9%?!<%5!<%P(B}
{$B%Q%C%1!<%8$N%9%]%s%5!<(B(mentor)$B$rC5$9%5!<%P(B}
{A}
{$B8E$$%5!<%P$O%G%#%9%/>c32Ey$,$"$C$?$N$G!"<wL?$HH=CG$5$l!"%G!<%?$,B;<:$9$kA0$K?7$7$$%5!<%P$K0\9T$5$l$^$7$?!#(Bftp-master.debian.org$B$O(BDebian$B$N(B official package $B%j%]%8%H%j$G$9!#%Q%C%1!<%8$N%9%]%s%5!<(B(mentor)$B$rC5$9$N$O(Bmentors.debian.net$B!#(B }

\santaku
{pbuilder$B$K(Bclang support$B$,DI2C$5$l$^$7$?!#C/$,=q$$$?%Q%C%A$G$7$g$&$+!)(B}
{Sylvestre Ledru}
{Junichi Uekawa}
{Hideki Yamane}
{C}
{Debian$B$N(BClang$B%5%]!<%H$OCe!9$H?J$s$G$$$^$9!#(B}

\santaku
{DPN - 2013$BG/(B3$B7n(B4$BF|9f$K<h$j>e$2$i$l$?F|K\$N%$%Y%s%H$O(B}
{Open Source Conference 2013 Tokyo/Spring}
{Open Source Conference 2013 Hamamatu}
{Open Source Conference 2013 Tokushima}
{A}
{\url{http://henrich-on-debian.blogspot.jp/2013/02/open-source-conference-2013-tokyospring.html} $B>\:Y$O8e$[$I!#(B}


\end{multicols}


% % (query-replace-regexp "<.*?>" "")
% % (query-replace-regexp "^[    ]\+" "")

%-------------------------------------------------------------------------------
\dancersection{DebianにおけるGo}{上川純一}
%-------------------------------------------------------------------------------

Go 言語をDebianで使ってみましょう。

まず入門するにはオンラインでのチュートリアルが充実しているので最初はロー
カルにインストールすることすら必要ないでしょう\cite{golang:tourofgo}
\cite{golang:tourofgoja}。しかし次第に本格的に使い始めようとすると手元に
インストールしたくなるはずです。

\subsection{DebianユーザとしてGoでプログラムを書いてみる}

Go のインストール手順\cite{golangorg-install}に従って公式パッケージをダ
ウンロードしてTarを展開すると最新のバイナリがインストールできます。しか
しあとからアンインストールするときとかに面倒ですしアップデートも自動では
ありません。ここはDebianパッケージが使いたいなぁと思います。

Debianの公式パッケージはBusterでは2021年4月現在1.11ですが、Buster
backports にはもう少し新しい1.14があります。まず /etc/apt/sources.listに
buster-backportsを追加しましょう。\footnote{時期安定版予定の bullseye だ
と1.15になっているようで、sidには1.16パッケージがあるようです。}

\begin{commandline}
deb https://deb.debian.org/debian buster-backports main
\end{commandline}

そしておもむろにインストール

\begin{commandline}
$ sudo apt update
$ sudo apt install -t buster-backports golang
$ sudo apt install -t buster-backports golang-mode  # もしEmacsユーザなら
\end{commandline}
%$ -- for emacs

Goの追加のパッケージはDebianパッケージではなくGoの仕組みでインストールし
ましょう。Debianパッケージはビルド時の依存関係を満たすためだけにパッケー
ジされているようです。理由は２つあり、

\begin{itemize}
 \item ユーザがgo getを使うことでgo build などに必要な環境をユーザが設定することになる。
 \item ユーザ権限では書き込めないのと、.gitディレクトリが存在しないため
       通常のgo get -u でのアップデートなどが不可能なためです。
\end{itemize}

以上の理由によりDebianパッケージでは\texttt{/usr/share/gocode/src/}にソー
スがインストールされるのですが特に指定しないとGoのコンパイル時に利用して
くれません。

たとえばsqlite3を利用したいなと思ったらgo getコマンドを利用してモジュー
ルを \verb!~/go! 以下にインストールします。 \footnote{Debianパッケージだ
とgolang-github-mattn-go-sqlite3-devがあるのだがそれはパッケージのビルド
用に用意されている。}

\begin{commandline}
$ go get github.com/mattn/go-sqlite3
\end{commandline}

すると以下のようなプログラム(db.go)がコンパイルできるようになります。

\begin{commandline}
import (
	"database/sql"
	_ "github.com/mattn/go-sqlite3"
	"log"
)

func main() {
	db, err := sql.Open("sqlite3", "./ac.db")
	if err != nil {
		log.Fatal(err)
	}
	defer db.Close()
}

\end{commandline}

コンパイルはgo buildで実行します。ただビルドするだけじゃなくてそのまま実
行までしたければgo runで実行できます。

\begin{commandline}
$ go build db.go
$ ./db

$ go run db.go
\end{commandline}
%$ -- for emacs

\subsubsection{GolangのModuleシステム}

Golangの標準のモジュール管理システムはGOPATHが従来利用されていたのですが、
1.11からGo Modules\cite{golang:go-modules}が提供され、1.14からデフォルト
に切り替わりました。

go.modファイルにモジュールメタデータが管理されています。\footnote{以前は
ディレクトリ構成とGOPATH環境変数で指定されたディレクトリで管理していまし
た}。ホームディレクトリ以下のgo/pkg以下にモジュールがインストールされる
ようです。\footnote{以前はホームディレクトリ以下の\texttt{go/src}}

Go Moduleを使うGoのコードを新規に開発するには、モジュールのディレクトリ
でgo mod initでgo.modを作成します。

\begin{commandline}
$ go mod init example.com/your/packagename
$ cat go.mod
module example.com/your/packagename

go 1.15

$ go test  # カレントディレクトリのモジュールのテストを走らせる
\end{commandline}
%$

サブディレクトリのパッケージのときにimport 文で
example.com/your/packagename/subdirのように指定するとローカルのサブディ
レクトリを利用してくれます。

それ以外の名前を指定したらgo build時にダウンロードしてきて、go.modに情報
を追記してくれるようです。

\subsubsection{DockerユーザとしてのGo}

DockerでDebianでGoを使う場合にはDockerhub のGolangのDockerイメージ
\cite{docker:go}を使うのが効率よくお手軽でしょう。

golangという名前でいくつかのバージョンが提供されておりデフォルトはBuster
イメージのようです。golang:latest は２０２１年４月現在では
golang:1.16.3-buster になっておりdocker pull やDockerfileにgolangと指定
するとDebianの上にGolangをインストールしたイメージが入るようです。

たとえばこんなDockerfileを書いてCIに利用しています。

\begin{commandline}
FROM golang
COPY . .
RUN make
\end{commandline}


\subsection{Debian開発者としてのGoパッケージ}

DebianにGoでできたソフトウェアをパッケージとしていれたいときには追加で作
業が必要になります。いちユーザとして利用していた場合との差異としてはポリ
シー準拠のためにDebianで依存関係すべて含めてビルドできるようにする必要が
あるということです。つまり依存関係にあったものをgo getでダウンロードして
くるのではなくDebianパッケージを利用します。

Debian にあるGoパッケージ作成はpkg-goチームが中心となって行っているよう
です\cite{debian-go-packaging}、全てではないですが大半のパッケージはこの
チームがやっているようです。

\texttt{dh-golang}を使ってDebianパッケージのビルドは行うようです。
\texttt{dh-make-golang}コマンドを使うとDebianパッケージの雛形を作成して
くれるようです。

\begin{commandline}
$ dh-make-golang make github.com/ncabatoff/fakescraper
2021/04/10 01:46:49 Starting "dh-make-golang v0.4.0 linux/amd64"
2021/04/10 01:46:49 Downloading "github.com/ncabatoff/fakescraper/..."
2021/04/10 01:46:50 Determining upstream version number
2021/04/10 01:46:50 Package version is "0.0~git20201102.4b37ba6"
2021/04/10 01:46:50 Determining dependencies
2021/04/10 01:46:50 Generating temp tarball as "/tmp/dh-make-golang230078837"
2021/04/10 01:46:50 Moving tempfile to "golang-github-ncabatoff-fakescraper_0.0~git20201102.4b37ba6.orig.tar.xz"
hint: Using 'master' as the name for the initial branch. This default branch name
hint: is subject to change. To configure the initial branch name to use in all
hint: of your new repositories, which will suppress this warning, call:
hint: 
hint: 	git config --global init.defaultBranch <name>
hint: 
hint: Names commonly chosen instead of 'master' are 'main', 'trunk' and
hint: 'development'. The just-created branch can be renamed via this command:
hint: 
hint: 	git branch -m <name>
2021/04/10 01:46:50 Adding remote "origin" with URL "git@salsa.debian.org:go-team/packages/golang-github-ncabatoff-fakescraper.git"
2021/04/10 01:46:50 Adding remote "github" with URL "https://github.com/ncabatoff/fakescraper"
2021/04/10 01:46:50 Running "git fetch github"
remote: Enumerating objects: 1, done.        
remote: Counting objects: 100% (1/1), done.        
remote: Total 14 (delta 0), reused 0 (delta 0), pack-reused 13        
Unpacking objects: 100% (14/14), 3.26 KiB | 238.00 KiB/s, done.
From https://github.com/ncabatoff/fakescraper
 * [new branch]      master     -> github/master
2021/04/10 01:46:52 Could not determine license for "github.com/ncabatoff/fakescraper": GET https://api.github.com/repos/ncabatoff/fakescraper/license: 404 Not Found []
2021/04/10 01:46:52 Setting debian/watch to track git HEAD
2021/04/10 01:46:52 Could not determine license for "github.com/ncabatoff/fakescraper": GET https://api.github.com/repos/ncabatoff/fakescraper/license: 404 Not Found []
2021/04/10 01:46:52 Done!

Packaging successfully created in /tmp/kk/golang-github-ncabatoff-fakescraper
    Source: golang-github-ncabatoff-fakescraper
    Binary: golang-github-ncabatoff-fakescraper-dev

Resolve all TODOs in itp-golang-github-ncabatoff-fakescraper.txt, then email it out:
    /usr/sbin/sendmail -t < itp-golang-github-ncabatoff-fakescraper.txt

Resolve all the TODOs in debian/, find them using:
    grep -r TODO debian

To build the package, commit the packaging and use gbp buildpackage:
    git add debian && git commit -a -m 'Initial packaging'
    gbp buildpackage --git-pbuilder

To create the packaging git repository on salsa, use:
    dh-make-golang create-salsa-project golang-github-ncabatoff-fakescraper

Once you are happy with your packaging, push it to salsa using:
    gbp push

NOTE: Full upstream git history has been included as per pkg-go team's
      new workflow.  This feature is new and somewhat experimental,
      and all feedback are welcome!
      (For old behavior, use --upstream-git-history=false)

The upstream git history is being tracked with the remote named "github".
To upgrade to the latest upstream version, you may use something like:
    git fetch github          # note the latest tag or commit-ish
    uscan --report-status     # check we get the same tag or commit-ish
    gbp import-orig --sign-tags --uscan --upstream-vcs-tag=<commit-ish>
 
\end{commandline}

\begin{thebibliography}{1}
 \bibitem{golangorg} ``The Go Programming Language'' \url{https://golang.org/}
 \bibitem{golangorg-doc} ``Documentation -- The Go Programming Language'', \url{https://golang.org/doc}
 \bibitem{golangorg-install} \url{https://golang.org/doc/install}
 \bibitem{golang:tourofgo} ``A Tour of Go'' \url{https://tour.golang.org/welcome/1}
 \bibitem{golang:tourofgoja} \url{https://go-tour-jp.appspot.com/welcome/1}
 \bibitem{debian-go-packaging} ``Debian Go Packaging'' \url{https://go-team.pages.debian.net/packaging.html}
 \bibitem{golang:go-modules} ```Go modules''' \url{https://github.com/golang/go/wiki/Modules}
 \bibitem{docker:go} ``golang Docker Official Images'' \url{https://hub.docker.com/_/golang}
\end{thebibliography}

%\dancersection{title}{username}


% 冊子にするために、4の倍数にする必要がある。
% そのための調整
\dancersection{メモ}{}
\mbox{}\newpage
\mbox{}\newpage

\vspace*{15cm}
\hrule
\vspace{2mm}
\includegraphics[width=2cm]{image-assets/openlogo-nd.eps}
\noindent \Large \bf Debian 勉強会資料\\
\noindent \normalfont \debmtgyear{}年\debmtgmonth{}月\debmtgdate{}日 \hspace{5mm}  初版第1刷発行\\
\noindent \normalfont 東京エリア Debian 勉強会 （編集・印刷・発行）\\
\hrule
\end{document}
